% =============================================================
% control_single.tex — Контрольная работа, один вариант, A4 portrait
% Компилятор: XeLaTeX
% =============================================================

\documentclass[12pt]{article}

% Сначала задаём портретную геометрию — ДО загрузки стиля,
% чтобы \geometry в .sty не перезаписал наши настройки.
\usepackage[a4paper, left=20mm, right=15mm, top=15mm, bottom=15mm]{geometry}

% Стиль — загружаем с флагом nogeom чтобы он не трогал геометрию
% Реализация: просто переопределим \geometry внутри .sty как no-op
% после того как он уже применился. Используем другой подход:
% грузим стиль, потом восстанавливаем геометрию.
\usepackage{../styles/mathworksheet}

% Восстанавливаем портретную геометрию после того как стиль
% задал landscape:
\newgeometry{a4paper, left=20mm, right=15mm, top=15mm, bottom=15mm,
             noheadfoot}

% ============================================================
% ПАРАМЕТРЫ РАБОТЫ
% ============================================================
\newcommand{\worktitle}{Контрольная работа №\,3}
\newcommand{\worktopic}{Линейные уравнения и их системы}
\newcommand{\workvariant}{2}
% ============================================================

\begin{document}

\settaskwidth{165mm}

% ============================================================
% ШАПКА
% ============================================================
\ctheader{\worktitle}{\worktopic}{\workvariant}

% ============================================================
% ЧАСТЬ 1. ВЫЧИСЛИ
% ============================================================
\ctsectionnum{1}{Вычисли}

\calcrowc{1}{\tfrac{3}{4} + \tfrac{5}{6}}
\calcrowc{2}{2{,}4 \cdot 1{,}5 - 0{,}6}
\calcrowc{3}{\tfrac{7}{12} - \tfrac{1}{4} + \tfrac{1}{3}}

% ============================================================
% ЧАСТЬ 2. РЕШИ УРАВНЕНИЯ
% ============================================================
\ctsectionnum{2}{Реши уравнения}

\taskrow{4}{$5x - 8 = 3x + 6$}
\vspace{10mm}

\taskrow{5}{$\dfrac{2x+1}{3} = \dfrac{x-2}{2}$}
\vspace{10mm}

\taskrow{6}{$3(2x - 5) - 2(x + 4) = 0$}
\vspace{10mm}

% ============================================================
% ЧАСТЬ 3. ЗАДАЧА
% ============================================================
\ctsectionnum{3}{Задача}

\ctproblem{7}{%
  Два поезда вышли одновременно навстречу друг другу
  из двух городов, расстояние между которыми равно $450$~км.
  Скорость первого поезда на $10$~км/ч больше скорости второго.
  Через $3$~часа они встретились. Найдите скорость каждого поезда.%
}{52}

\end{document}

% ============================================================
% СПРАВОЧНИК КОМАНД
% ============================================================
% Шапка:       \ctheader{название}{тема}{вариант}
% Секции:      \ctsectionnum{N}{название}
%              \ctsection{текст}
% Вычисления:  \calcrowc{N}{выр}   \calcrow{N}{выр}
% Задания:     \taskrow{N}{текст}  \taskrowans{N}{текст}
% Задача:      \ctproblem{N}{условие}{высота_мм}
% Прочее:      \tasksep  \ansline[ширина]  \mf{формула}
%
% КОМПИЛЯЦИЯ:  xelatex control_single.tex  (два прохода)
% ============================================================